\documentclass{article}
\usepackage{amsmath}
\usepackage{graphicx}
\graphicspath{ {images/} }
\usepackage[spanish]{babel}
\usepackage{bbm}
\usepackage{amsthm}
\usepackage{authblk}
\usepackage{hyperref}
\usepackage[utf8]{inputenc}
\usepackage{mathrsfs}
\usepackage{amsfonts}
\usepackage{amssymb}
\usepackage{enumerate}
\usepackage[left=3cm,right=2cm,top=2cm,bottom=2cm]{geometry}
\usepackage{epsfig}
 \renewcommand{\baselinestretch}{0.93}
 \thispagestyle{empty}
 \pagestyle{empty}
\setlength{\textheight}{53pc} \setlength{\textwidth} {36pc}
\setlength{\topmargin}{-4mm}
\usepackage{multicol}
\newcommand*{\email}[1]{%
    \normalsize\href{mailto:#1}{#1}\par
    }
\setlength{\columnsep}{1cm}
\newcommand{\N}{\mathbb{N}}
\newcommand{\calM}{\cal{M}}
\newcommand{\calP}{\cal{P}}
\newcommand{\F}{\mathbb{F}}
\newcommand{\R}{\mathbb{R}}
\newcommand{\Z}{\mathbb{Z}}
\newcommand{\Q}{\mathbb{Q}}
\newcommand{\C}{\mathbb{C}}
\newcommand{\bbH}{\mathbb{H}}


\theoremstyle{definition}
\newtheorem{ejer}{Ejercicio}
\author{Marcos Morales}
\affil{Universidad Finis Terrae}
\title{Semana 4\\}



	


\begin{document}


\begin{picture}(400, 75)
   \put(-40,15){\includegraphics[width=5cm]{UFT.png}}


\end{picture}


\begin{center}
\huge{Laplace y Aplicaciones}
\end{center}

\begin{center}
\Large{Marcos Morales, Camila Muñoz}
\end{center}

\begin{center}
	\large{Ecuaciones Diferenciales Ordinarias (EDO)}
\end{center}
\begin{center}
\setlength{\parindent}{0cm}\large{Clases centradas para capacitar a los estudiantes de ingeniería en Ecuaciones diferenciales ordinarias. \\}
\end{center}


\vspace{1cm}

\begin{center}
	\large{Contenidos:}
\end{center}
\phantom{abefwfewefwfefffwi} - Ecuación de Cauchy-Euler\\
\phantom{abefwfewefwfefffwi} - Transformada de Laplace\\
\phantom{abefwfewefwfefffwi} - Transformada Inversa\\
\phantom{abefwfewefwfefffwi} - Transformada de una derivada\\





\vspace{6cm}


\begin{center}
\today
\end{center}
\begin{center}
Santiago, Chile
\end{center}


\pagestyle{empty}


\setcounter{page}{1}
\pagestyle{plain}
\setlength{\parindent}{0cm}





\newpage
\section{Ecuación de Cauchy-Euler}

La ecuación de Cauchy-Euler es de la forma


\begin{center}
$\displaystyle a_nx^ny^{(n)}+a_{n-1}x^{n-1}y^{(n-1)}+... +a_1xy'+a_0y = 0$
\end{center}

En concreto vamos a analizar el caso cuando $n=2$, es decir

\begin{center}
$\displaystyle ax^2y''+bxy'+cy =0$
\end{center}


En este caso se asume que la solución es de la forma $y=x^m$, entonces $y'=mx^{m-1}$ y $y''=m(m-1)x^{m-2}$. Reemplazando en la ecuación diferencial inicial obtenemos


\begin{center}
$\displaystyle ax^2m(m-1)x^{m-1}+bxmx^{m-1}+cx^m =0$
\end{center}

Luego


\begin{center}
$\displaystyle am(m-1)+bm+c =0$
\end{center}

\begin{center}
$\displaystyle am^2+(b-a)m+c =0$
\end{center}

Esta es una ecuación cuadrática con solución  es


\begin{center}
$\displaystyle m = \frac{a-b \pm \sqrt{(b-a)^2-4ac}}{2a}$
\end{center}

Existen 3 casos. \textbf{Caso 1: } \\


$(b-a)^2-4ac >0$, en este caso hay dos soluciones



\begin{center}
$\displaystyle m_1 = \frac{a-b + \sqrt{(b-a)^2-4ac}}{2a}$
\end{center}


\begin{center}
$\displaystyle m_2 = \frac{a-b - \sqrt{(b-a)^2-4ac}}{2a}$
\end{center}

Y por lo tanto las soluciones de la ecuación diferencial son


\begin{center}
$y_1=x^{m_1}$ y $y_2=x^{m_2}$
\end{center}

Por lo tanto la solución general es

\begin{center}
$y= c_1x^{m_1}+c_2x^{m_2}$
\end{center}

\textbf{Caso 2: } \\


$(b-a)^2-4ac = 0$, en este caso hay sólo una solución



\begin{center}
$\displaystyle m_1 = \frac{a-b}{2a}$
\end{center}


Luego $y_1= x^{m_1}$ es solución de la ecuación diferencial. Para encontrar la segunda solución usamos la solución

\begin{center}
$\displaystyle y_2 = y_1 \int \frac{e^{\int -\frac{b}{ax} dx}}{y_1^2} dx$
\end{center}


Resolviendo nos queda

\begin{center}
$\displaystyle y_2 = x^{m_1}\ln(x)$
\end{center}

Por lo tanto la solución general es

\begin{center}
$y= c_1x^{m_1}+c_2x^{m_1}\ln(x)$
\end{center}


\textbf{Caso 2: } \\


$(b-a)^2-4ac < 0$, en este caso hay dos soluciones imaginarias



\begin{center}
$\displaystyle m_1 = \alpha+ \beta i$
\end{center}

\begin{center}
$\displaystyle m_2 = \alpha - \beta i$
\end{center}


Realizando los mismos pasos que se usaron en las ecuaciones de coeficientes constantes, obtenemos las soluciones

\begin{center}
$y_1 =  x^{\alpha}\cos(\beta \ln(x))$ y  $ y_2 =x^{\alpha}\sen(\beta \ln(x))$
\end{center}


Por lo tanto la solución general será de la forma 


\begin{center}
$y = c_1 x^{\alpha}\cos(\beta \ln(x)) +c_2 x^{\alpha}\sen(\beta \ln(x)) $
\end{center}


Con esto hemos resuleto todos los casos posibles. \\





\textbf{Ejemplo: }Resolver la ecuación diferencial $x^2y''-3xy'+3y=2x^4e^x$. \\


\textit{Respuesta:} Primero resolvemos la parte homogenea $x^2y''-3xy'+3y =0$, usando $y=x^m$ obtenemos

\begin{center}
$m(m-1)-3m+3=0$
\end{center}


\begin{center}
$m^2-4m+3=0$
\end{center}

Las soluciones son $m_1= 1$ y $m_2=3$ , por lo tanto la solución de la homogenea es $y_h = c_1 x+c_2x^3$. Ahora bien para calcular la solución particular $y_p$ usaremos variación de parámetros y supondremos que la solución es de la forma $y_p =c_1(x)x +c_2(x)x^3$. Primero calcularemos el Wronskiano

\begin{center}
$W(y_1,y_2) = x(3x^2)-x^3 = 2x^3$
\end{center}

Segundo, dejamos la ecuación diferencial inicial en la forma $y''+p(x)y'+q(x)y=f(x)$, luego


\begin{center}
$\displaystyle y'' -\frac{3}{x} y'+\frac{3}{x^2} y=2x^2e^x$
\end{center}

Por lo tanto $f(x) = 2x^2e^x$. Por último calculamos $c_1(x)$ y $c_2(x)$

\begin{align*}
\displaystyle c_1(x) &= \int \frac{-y_2 f}{W(y_1,y_2)} dx \\
&= \int \frac{-x^3 2x^2e^x}{2x^3} dx \\
&= -\int x^2e^x dx \\
\end{align*}

Para resolver esta integral usamos integrales por partes $u=x^2$ y $dv=e^xdx$, luego $du=2xdx$ y $v=e^x$

\begin{align*}
\displaystyle  \int x^2e^x dx  &= x^2e^x -\int 2xe^x dx \\
&= x^2e^x -2\int xe^x dx
\end{align*}


Ahora para integrar $xe^x$ usamos nuevamente integrales por partes $u= x$ y $dv=e^x$, entonces $du=dx$ y $v=e^x$, luego


\begin{align*}
\displaystyle  \int xe^x dx  &= xe^x -\int e^x dx \\
&= xe^x -e^x
\end{align*}

Por lo tanto

\begin{align*}
\displaystyle  \int x^2e^x dx  &= x^2e^x -2(xe^x-e^x) 
\end{align*}


\begin{align*}
\displaystyle  \int x^2e^x dx  &= x^2e^x -2xe^x+2e^x 
\end{align*}


Por lo tanto 

\begin{align*}
\displaystyle c_1(x) &= 2xe^x-x^2e^x -2e^x  \\
\end{align*}


Por otra parte 


\begin{align*}
\displaystyle c_2(x) &= \int \frac{y_1 f}{W(y_1,y_2)} dx \\
&= \int \frac{x 2x^2e^x}{2x^3} dx \\
&= \int e^x dx \\
&= e^x
\end{align*}


Por lo tanto la solución general es

\begin{center}
$y= c_1 x+c_2x^3 + x(2xe^x-x^2e^x -2e^x)+ x^3e^x $
\end{center}

\newpage





\section{Transformada de Laplace}

Sea $f: [0, \infty) \to \mathbb{R}$. Entonces la integral

\begin{center}
$\displaystyle \int_0^\infty f(t)e^{-st} dt$
\end{center}

Se denomina la transformada de Laplace de $f$ y se denota por

\begin{center}
$\mathcal{L}[f](s)$, $F(s)$ o bien $\hat{f}(s)$
\end{center}

Siempre que la integral exista, nosotros usaremos la notación $\mathcal{L}[f](s)$. \\


La transformada de Laplace es un funcional lineal, es decir, toma una función $f$ y la manda a otra función $\mathcal{L}[f]$. Notar que $\mathcal{L}[f](s)$ es una función que depende de $s$, no es un número. \\


Vamos a calcular directamente algunas transformadas de Laplace, por ejemplo, si tomamos $f(t) = 1$, entonces


\begin{align*}
\displaystyle \mathcal{L}[1](s) &= \int_0^\infty 1 \cdot e^{-st} dt \\
&= \int_0^\infty e^{-st} dt \\
&= \lim_{A \to \infty} \int_0^A e^{-st} dt \\
&= \lim_{A \to \infty} -\frac{1}{s}\int_0^A e^{u} du \\
&= \lim_{A \to \infty} -\frac{1}{s}e^{-st} \bigg|_0^A \\
&= \lim_{A \to \infty} -\frac{1}{s}e^{-sA} + \frac{1}{s}e^{s \cdot 0}\\
&= \frac{1}{s}\
\end{align*}

Tomamos ahora $f(t) = t$, entonces

\begin{align*}
\displaystyle \mathcal{L}[t](s) &= \int_0^\infty 1 \cdot te^{-st} dt \\
&= \lim_{A \to \infty} \int_0^A te^{-st} dt \\
\end{align*}

Vamos a calcular $\displaystyle \int te^{-st}dt$, tomamos $u=t$ y $dv=e^{-st}dt$, entonces $du=dt$ y $v=\displaystyle -\frac{1}{s}e^{-st}$, luego

\begin{align*}
\displaystyle \int te^{-st}dt &= -\frac{te^{-st}}{s}- \int -\frac{1}{s}e^{-st} dt \\
&= -\frac{te^{-st}}{s}- \frac{1}{s^2}e^{-st} \\
\end{align*}


Por lo tanto

\begin{align*}
\displaystyle \mathcal{L}[t](s) &= \lim_{A \to \infty} \int_0^A te^{-st} dt \\
&= \lim_{A \to \infty} \left( -\frac{te^{-st}}{s}- \frac{1}{s^2}e^{-st}  \right)\bigg|_0^A  \\
&= \lim_{A \to \infty} \left( -\frac{Ae^{-sA}}{s}- \frac{1}{s^2}e^{-sA}  -  \left( -\frac{0 \cdot e^{-s \cdot 0}}{s}- \frac{1}{s^2}e^{-s\cdot 0} \right) \right) \\
&= \frac{1}{s^2}
\end{align*}



Tomamos ahora $f(t) = \sen(at)$, entonces

\begin{align*}
\displaystyle \mathcal{L}[\sen(at)](s) &= \int_0^\infty 1 \cdot \sen(at)e^{-st} dt \\
&= \lim_{A \to \infty} \int_0^A \sen(at)e^{-st}  dt \\
\end{align*}

Vamos a calcular $\displaystyle \int \sen(at)e^{-st}dt$, tomamos $u=\sen(at)$ y $dv=e^{-st}dt$, entonces $du=a\cos(at)dt$ y $v=\displaystyle -\frac{1}{s}e^{-st}$, luego

\begin{align*}
\displaystyle \int \sen(t)e^{-st}dt &= -\frac{\sen(at)e^{-st}}{s}- \int -\frac{1}{s}e^{-st} a\cos(at)dt \\
&= -\frac{\sen(at)e^{-st}}{s}+ \frac{a}{s} \int \cos(at)e^{-st} dt\\
\end{align*}

Vamos a calcular $\displaystyle \int \cos(at)e^{-st}dt$, tomamos $u=\cos(at)$ y $dv=e^{-st}dt$, entonces $du=-a\sen(at)dt$ y $v=\displaystyle -\frac{1}{s}e^{-st}$, luego



\begin{align*}
\displaystyle \int \cos(at)e^{-st}dt &= -\frac{\cos(at)e^{-st}}{s}- \int (-\sen(at))\frac{-1}{s} e^{-st} adt \\
&= -\frac{\cos(at)e^{-st}}{s}- \frac{a}{s} \int \sen(at)e^{-st} dt\\
\end{align*}

Reemplazando

\begin{align*}
\displaystyle \int \sen(t)e^{-st}dt &= -\frac{\cos(at)e^{-st}}{s}- \int -\frac{1}{s}e^{-st} a\cos(at)dt \\
&= -\frac{\sen(at)e^{-st}}{s}+ \frac{a}{s} \left( -\frac{\cos(at)e^{-st}}{s}- \frac{a}{s} \int \sen(at)e^{-st} dt \right)  \\
&= -\frac{\sen(at)e^{-st}}{s}+  -\frac{a\cos(at)e^{-st}}{s^2}- \frac{a^2}{s^2} \int \sen(at)e^{-st} dt \\
\end{align*}

Ahora sumando $\displaystyle \frac{a^2}{s^2} \int \sen(at)e^{-st} dt $ a ambos lados obtenemos

\begin{align*}
\displaystyle \int \sen(t)e^{-st}dt + \frac{a^2}{s^2} \int \sen(at)e^{-st} dt &= -\frac{\sen(at)e^{-st}}{s} -\frac{a\cos(at)e^{-st}}{s^2} \\
\end{align*}

O bien 


\begin{align*}
\displaystyle  \int \sen(at)e^{-st} dt  \left( 1+ \frac{a^2}{s^2}\right) &= -\frac{\sen(at)e^{-st}}{s} -\frac{a\cos(at)e^{-st}}{s^2} \\
\end{align*}

\begin{align*}
\displaystyle  \int \sen(at)e^{-st} dt  \left( \frac{a^2+s^2}{s^2}\right) &= -\frac{\sen(at)e^{-st}}{s}  -\frac{a\cos(at)e^{-st}}{s^2} \\
\end{align*}


\begin{align*}
\displaystyle  \int \sen(at)e^{-st} dt &= \left( -\frac{\sen(at)e^{-st}}{s} -\frac{a\cos(at)e^{-st}}{s^2} \right) \frac{s^2}{a^2+s^2}\\
\end{align*}

Por lo tanto 


\begin{align*}
\displaystyle \mathcal{L}[\sen(t)](s) &= \lim_{A \to \infty} \int_0^A \sen(t)e^{-st}  dt \\
&= \lim_{A \to \infty} \left( \left( -\frac{\sen(at)e^{-st}}{s} -\frac{a\cos(at)e^{-st}}{s^2} \right) \frac{s^2}{a^2+s^2}\right) \bigg|_0^A \\
&= \frac{a}{a^2+s^2}
\end{align*}


En resumen

\begin{align*}
\displaystyle \mathcal{L}[\sen(at)](s) &= \frac{a}{a^2+s^2}
\end{align*}


por último, vamos a calcular la transformada de Laplace de $e^{\alpha t}$, tenemos uqe

\begin{align*}
\displaystyle \mathcal{L}[e^{\alpha t}](s) &= \int_0^\infty e^{\alpha t} e^{-st} dt \\
&= \lim_{A \to \infty} \int_0^A e^{t(\alpha-s)} dt \\
&= \lim_{A \to \infty} \left( \frac{1}{\alpha -s} e^{t(\alpha-s)} \right)\bigg|_0^A  \\
&= \frac{1}{s-\alpha}
\end{align*}

Esto último es cierto cuanod $s >\alpha$.

\newpage

A continuación dejamos una lista de transformadas de Laplace


\begin{center}
$\displaystyle \mathcal{L}[1](s) = \frac{1}{s}$, $\forall s>0$
\end{center}

\begin{center}
$\displaystyle \mathcal{L}[t^n](s) = \frac{n!}{s^{n+1}}$, $\forall s>0$ y $n \in \mathbb{N}$
\end{center}

\begin{center}
$\displaystyle \mathcal{L}[e^{\alpha t}](s) = \frac{1}{s-\alpha}$, $\forall s>\alpha$
\end{center}

\begin{align*}
\displaystyle \mathcal{L}[\sen(at)](s) &= \frac{a}{a^2+s^2}, \forall s>0
\end{align*}

\begin{align*}
\displaystyle \mathcal{L}[\cos(at)](s) &= \frac{s}{a^2+s^2} \forall s>0
\end{align*}


La siguiente propiedad es fundamental para calcular transformadas de Laplace\\

\textbf{Teorema: } La transformada de Laplace es una transformación lineal. \\

Esto significa que la transformada de laplace separa sumas y saca afuera las contantes. \\

\textbf{Ejemplo: } Calcular la transformada de Laplace de $f(t)=1+4t^4+\cos(2t)$. \\

\textit{Respuesta: } Tenemos que

\begin{align*}
\displaystyle \mathcal{L}[1+4t^4+\cos(2t)](s) &= \mathcal{L}[1](s)+ 4 \mathcal{L}[t^4](s)+\mathcal{L}[\cos(2t)](s) \\
&= \frac{1}{s} + 4 \frac{4!}{s^5} + \frac{s}{2^2+s^2}\\
&= \frac{1}{s} + \frac{96}{s^5} + \frac{s}{4+s^2}\\
&= \frac{s^4(4+s^2)+ 96(4+s^2) + s \cdot s^5}{s^5(4+s^2)}\\
&= \frac{2s^6 +4s^4+96s^2+384}{s^5(4+s^2)}\\
\end{align*}


\newpage



\section{Transformada Inversa}


\textbf{Definición: }Una función $f$ se dice de orden exponencial si existe $\alpha \in \mathbb{R}$ y $M >0$ tal que	

\begin{center}
$\displaystyle |f(t)| < Me^{\alpha t}$  $\forall t>0$
\end{center}


\textbf{Ejemplo: } La función $\cos(x)$ es de orden exponencial pues $|\cos(x)| \leq 1 = 1 \cdot e^{0 t}$, por otra parte la función exponencial $e^{ax}$ también es de orden exponencial pues $|e^{ax}| =e^{ax} \leq 1 \cdot e^{ax}$. Se tiene además que cualquier polinomio es de orden exponencial. \\



\textbf{Definición: } Diremos que una función es continua por trozos si tiene a lo más un número finito de discontinuidades. \\

Con estas definiciones tenemos el siguiente teorema. \\

\textbf{Teorema: } Sea $f$ función continua a trozos en $(0,\infty)$ y de orden exponencial, entonces existe $\mathcal{L}[f]$. Además, si $g$ es otra función continua a trozos de orden exponecial tal que $\mathcal{L}[f]=\mathcal{L}[g]$, entonces se tiene que $f=g$ \\ \\

Este último teorema implica que $\mathcal{L}$ es un funcional lineal inyectivo y por lo tanto existe la inversa $\mathcal{L}^{-1}$ llamamos a esta función la \textbf{Transformada inversa de $f$}. \\

Los siguientes resultados son consecuencia directa de lo demostrado para la transformada de Laplace


\begin{center}
$\displaystyle \mathcal{L}^{-1}\left[ \frac{1}{s} \right] (t) = 1$
\end{center}


\begin{center}
$\displaystyle \mathcal{L}^{-1}\left[ \frac{1}{s^n} \right] (t) = \frac{t^{n-1}}{(n-1)!}$
\end{center}

\begin{center}
$\displaystyle \mathcal{L}^{-1}\left[ \frac{1}{s-a} \right] (t) = e^{at}$
\end{center}


\begin{center}
$\displaystyle \mathcal{L}^{-1}\left[ \frac{k}{s^2+k^2} \right] (t) = \sen(kt)$
\end{center}



\begin{center}
$\displaystyle \mathcal{L}^{-1}\left[ \frac{s}{s^2+k^2} \right] (t) = \cos(kt)$
\end{center}

\begin{center}
$\displaystyle \mathcal{L}^{-1}\left[ \frac{k}{s^2-k^2} \right] (t) = \sinh(kt)$
\end{center}

\begin{center}
$\displaystyle \mathcal{L}^{-1}\left[ \frac{s}{s^2-k^2} \right] (t) = \cosh(kt)$
\end{center}


\newpage



La transformada inversa es un transformación lineal al igual que la transformada de Laplace \\

\textbf{Ejemplo: }


\begin{center}
$\displaystyle \mathcal{L}^{-1}\left[ \frac{1}{s^3} + \frac{5}{s-3}\right](s)= \mathcal{L}^{-1}\left[ \frac{1}{s^3}\right](s) + 5\mathcal{L}^{-1} \left[ \frac{1}{s-3}\right](s) = \frac{t^2}{2} + 5e^{3t}$
\end{center}



Para calcular la transformada inversa de una fracción se descompone la fracción usando fracciones parciales \\


\textbf{Ejemplo 1: } Calcular $\displaystyle \mathcal{L}^{-1}\left[ \frac{s+1}{s^2(s+2)} \right] $. \\

\textit{Respuesta: } Tenemos que usar

\begin{align*}
\displaystyle  \frac{s+1}{s^2(s+2)} &= \frac{A}{s}+\frac{B}{s^2}+\frac{C}{s+2} \\
&= \frac{As(s+2)+ B(s+2)+Cs^2}{s^2(s+2)} \\
&= \frac{(A+C)s^2+(2A+B)s+2B}{s^2(s+2)} \\
\end{align*}

De donde se concluye que $A+C=0$, $2A+B=1$ y $2B=1$, luego $B=\frac{1}{2}$, $A=\frac{1}{4}$ y $C=- \frac{1}{4}$. Por lo tanto



\begin{align*}
\displaystyle  \frac{s+1}{s^2(s+2)} &= \frac{1}{4s}+\frac{1}{2s^2}-\frac{1}{4(s+2)} \\
\end{align*}

Por lo tanto



\begin{align*}
\displaystyle  \mathcal{L}^{-1}\left[ \frac{s+1}{s^2(s+2)} \right](t) &= \mathcal{L}^{-1}\left[ \frac{1}{4s}+\frac{1}{2s^2}-\frac{1}{4(s+2)} \right](t) \\
 &= \mathcal{L}^{-1}\left[ \frac{1}{4s}\right] + \mathcal{L}^{-1}\left[ \frac{1}{2s^2} \right]- \mathcal{L}^{-1}\left[ \frac{1}{4(s+2)} \right] \\
 &= \frac{1}{4}+ \frac{t}{2}-\frac{e^{-2t}}{4}
\end{align*}


\textbf{Ejemplo 2: } Calcular $\displaystyle \mathcal{L}^{-1}\left[ \frac{3s+4}{s^2+7} \right] $. \\

\textit{Respuesta: } En este caso no se usan fracciones parciales porque el denominador no se puede factorizar

\begin{align*}
\displaystyle  \frac{3s+4}{s^2+7} = \frac{3s}{s^2+7}  + \frac{4}{s^2+7} \\
\end{align*}

Luego, hay que calcular



\begin{center}
$\displaystyle \mathcal{L}^{-1}\left[ \frac{3s+4}{s^2+7} \right] = 3\displaystyle \mathcal{L}^{-1}\left[ \frac{s}{s^2+7} \right] + 4\displaystyle \mathcal{L}^{-1}\left[ \frac{1}{s^2+7} \right] $
\end{center}

Notemos que 

\begin{center}
$\displaystyle \mathcal{L}^{-1}\left[ \frac{k}{s^2+k^2} \right] (t) = \sen(kt)$
\end{center}



\begin{center}
$\displaystyle \mathcal{L}^{-1}\left[ \frac{s}{s^2+k^2} \right] (t) = \cos(kt)$
\end{center}

Por lo tanto tomando $k = \sqrt{7}$, tenemos que $k^2=7$, luego




\begin{center}
$\displaystyle \mathcal{L}^{-1}\left[ \frac{3s+4}{s^2+7} \right] = 3 \cos(\sqrt{7}t) + \frac{4}{\sqrt{7}}\sen(\sqrt{7}t) $
\end{center}


\newpage

\section{Transformada de una derivada}


\textbf{Teorema: }Supongamos que $f$ es continua en $(0,+\infty)$ y que además $f$ es de orden exponencial. Además supongamos que $f'$ es continua por tramos y también de orden exponencial. Entonces


\begin{center}
$\displaystyle \mathcal{L}[f'(t)](s) = s\mathcal{L}[f](s)- f(0)$
\end{center}

Más generalmente

\begin{center}
$\displaystyle \mathcal{L}[f^{(n)}(t)](s) = s^n\mathcal{L}[f](s)- s^{n-1}f(0)-s^{n-2}f'(0) - ... -f^{(n-1)}(0)$
\end{center}

En particular para $n=2$


\begin{center}
$\displaystyle \mathcal{L}[f''(t)](s) = s^2\mathcal{L}[f](s)- sf(0)-f'(0)$
\end{center}

Esto nos permite calcular ciertas transformadas de Laplace \\

\textbf{Ejemplo: } Caclular $\displaystyle \mathcal{L}\left[ \sen^2(at) \right] $. \\

notemos que $\sen^2(a \cdot 0) =0$, además

\begin{center}
$(\sen^2(at))' = 2a\cos(at)\sen(at) = a\cos(2at)$
\end{center}

Luego


\begin{center}
$ \displaystyle \mathcal{L}\left[ (\sen^2(at))'\right]  = a\mathcal{L}\left[ \cos(2at)\right] = \frac{2a^2}{s^2+4a^2} $
\end{center}

Por lo tanto 

\begin{center}
$ \displaystyle \mathcal{L}\left[\sen^2(at)\right]  = \frac{2a^2}{s(s^2+4a^2)} $
\end{center}



\section{Transformada de una integral definida}


Al igual que en la sección anterior, el siguiente resultado puede ser demostrado directamente usando la definición de la transformada de Laplace. Pero no vamos a demostrarlo y tampoco lo vamos a usar mucho por lo tanto no daremos ejemplos, tampoco haremos ejercicios\\

\textbf{Proposición: } Sea $f$ una función de orden exponencial, entonces

\begin{align*}
    \mathcal{L}\left[ \int_0^t f(u)du \right](s) = \frac{\mathcal{L}[f(t)](s)}{s}
\end{align*}













\end{document} 